\documentclass[12pt,a4paper]{article}

\usepackage[T2A]{fontenc}
\usepackage[utf8]{inputenc}
\usepackage[russian]{babel}
\usepackage{amsmath,amssymb}
\usepackage{geometry}
\usepackage{enumitem}
\geometry{margin=2cm}
\setlist[enumerate]{itemsep=0.6em}

\title{Тренажёр «Найди ошибку в решении»}
\author{Математика: вычисления, логарифмы, тригонометрия}
\date{}

\begin{document}
\maketitle

\section*{Как работать с пособием}
В каждом пункте дано \textbf{готовое решение} с одной или несколькими \textbf{преднамеренными ошибками}.
Задача ученика:
\begin{enumerate}
    \item найти место(а) ошибки;
    \item объяснить, почему это ошибка;
    \item записать корректное решение;
    \item выполнить самопроверку (подстановка, оценка знака, проверка ОДЗ, проверка периода и т.д.).
\end{enumerate}

\section{Разноуровневые вычислительные задачи}

\subsection*{Задание 1 (базовый уровень)}
Вычислить: $48:6\cdot 2-7$.

\textbf{Решение (с ошибкой):}
\[
48:6\cdot 2-7 = 48:12-7 = 4-7=-3.
\]

\subsection*{Задание 2 (средний уровень)}
Вычислить: $\dfrac{3}{4}-\dfrac{5}{8}+\dfrac{1}{2}$.

\textbf{Решение (с ошибкой):}
\[
\frac{3}{4}-\frac{5}{8}+\frac{1}{2}
= \frac{3-5+1}{8}
= -\frac{1}{8}.
\]

\subsection*{Задание 3 (повышенный уровень)}
Вычислить: $\left(2^{-3}+\dfrac{1}{4}\right)^{-1}$.

\textbf{Решение (с ошибкой):}
\[
\left(2^{-3}+\frac{1}{4}\right)^{-1}
= \left(-8+\frac{1}{4}\right)^{-1}
= \left(-\frac{31}{4}\right)^{-1}
= -\frac{31}{4}.
\]

\section{Логарифмические вычисления}

\subsection*{Задание 4}
Вычислить: $\log_2 32 + \log_2 \dfrac{1}{8}$.

\textbf{Решение (с ошибкой):}
\[
\log_2 32 + \log_2 \frac{1}{8}
= \log_2\left(32+\frac{1}{8}\right)
= \log_2\frac{257}{8}.
\]

\subsection*{Задание 5}
Вычислить: $\log_3 81 - \log_3 9$.

\textbf{Решение (с ошибкой):}
\[
\log_3 81 - \log_3 9 = \log_3(81-9)=\log_3 72.
\]

\subsection*{Задание 6}
Вычислить: $\log_{0.5}4$.

\textbf{Решение (с ошибкой):}
\[
\log_{0.5}4=2,
\]
так как $0.5^2=4$.

\section{Логарифмические уравнения}

\subsection*{Задание 7}
Решить уравнение: $\log_2(x-1)=3$.

\textbf{Решение (с ошибкой):}
\[
x-1=3^2=9 \Rightarrow x=10.
\]
Проверка не требуется.

\subsection*{Задание 8}
Решить уравнение: $\log_3 x + \log_3(x-2)=1$.

\textbf{Решение (с ошибкой):}
\[
\log_3 x + \log_3(x-2)=\log_3\bigl(x+x-2\bigr)=1,
\]
\[
\log_3(2x-2)=1 \Rightarrow 2x-2=3 \Rightarrow x=\frac{5}{2}.
\]

\subsection*{Задание 9}
Решить уравнение: $\log_5(x^2-4x+3)=\log_5(2x-2)$.

\textbf{Решение (с ошибкой):}
Из равенства логарифмов:
\[
x^2-4x+3=2x-2,
\]
\[
x^2-6x+5=0 \Rightarrow x=1 \text{ или } x=5.
\]
Ответ: $1;5$.

\section{Тригонометрические вычисления}

\subsection*{Задание 10}
Вычислить: $\sin\dfrac{\pi}{6}+\cos\dfrac{\pi}{3}$.

\textbf{Решение (с ошибкой):}
\[
\sin\frac{\pi}{6}+\cos\frac{\pi}{3}
= \frac{1}{2}+\frac{\sqrt3}{2}
= \frac{1+\sqrt3}{2}.
\]

\subsection*{Задание 11}
Вычислить: $\tan\dfrac{\pi}{4}\cdot\sin\dfrac{\pi}{2}$.

\textbf{Решение (с ошибкой):}
\[
\tan\frac{\pi}{4}\cdot\sin\frac{\pi}{2}=0\cdot1=0.
\]

\subsection*{Задание 12}
Вычислить: $\cos\left(-\dfrac{\pi}{3}\right)$.

\textbf{Решение (с ошибкой):}
\[
\cos\left(-\frac{\pi}{3}\right)=-\cos\frac{\pi}{3}=-\frac{1}{2}.
\]

\section{Тригонометрические уравнения}

\subsection*{Задание 13}
Решить уравнение: $\sin x=\dfrac{1}{2}$.

\textbf{Решение (с ошибкой):}
\[
x=\frac{\pi}{6}+2\pi n,\quad n\in\mathbb Z.
\]
Это все решения.

\subsection*{Задание 14}
Решить уравнение: $\cos x=-1$.

\textbf{Решение (с ошибкой):}
\[
x=\pi+\pi n,\quad n\in\mathbb Z.
\]

\subsection*{Задание 15}
Решить уравнение: $2\sin x\cos x=1$.

\textbf{Решение (с ошибкой):}
\[
2\sin x\cos x=\sin x,
\]
значит
\[
\sin x=1 \Rightarrow x=\frac{\pi}{2}+2\pi n,\quad n\in\mathbb Z.
\]

\section*{Блок самопроверки ученика}
После исправления каждого задания отметь:
\begin{itemize}
    \item Проверил(а) ли я область допустимых значений?
    \item Проверил(а) ли я свойства логарифмов перед применением формул?
    \item Проверил(а) ли я периодичность и полноту ответа в тригонометрии?
    \item Сделал(а) ли я подстановку корней в исходное уравнение?
\end{itemize}

\end{document}
